% Generated draft: requires mathematical and source review.
\documentclass[11pt]{article}
\usepackage[margin=25mm]{geometry}
\usepackage{fontspec}
\setmainfont{Latin Modern Roman}
\usepackage{amsmath,amssymb,mathtools,unicode-math}
\setmathfont{Latin Modern Math}
\usepackage{enumitem,needspace,xurl,hyperref}
\setlength{\parindent}{0pt}
\setlength{\parskip}{5pt}
\setlength{\emergencystretch}{3em}
\newcommand{\N}{\mathbb{N}}
\newcommand{\Z}{\mathbb{Z}}
\newcommand{\Q}{\mathbb{Q}}
\newcommand{\R}{\mathbb{R}}
\newcommand{\C}{\mathbb{C}}
\newcommand{\F}{\mathbb{F}}
\newcommand{\A}{\mathbb{A}}
\newcommand{\PP}{\mathbb{P}}
\newcommand{\E}{\mathbb{E}}
\newcommand{\eps}{\varepsilon}
\newcommand{\dd}{\,\mathrm{d}}
\begin{document}
\section*{UnsolvedMath ID: 800895}
\textbf{Source problem code:} catalogue entry
\subsection{Definitions and mathematical statement}
Let \(n\geq 1\) be an integer. A domain \(\Omega\subset\C^n\) is called hyperconvex if it admits a continuous plurisubharmonic exhaustion \(\rho:\Omega\to(-\infty,0)\) such that \(\{z\in\Omega:\rho(z)<c\}\) is relatively compact in \(\Omega\) for every \(c<0\). Here a plurisubharmonic function is an upper semicontinuous function whose restriction to each complex line is subharmonic or identically \(-\infty\).

For \(w\in\Omega\), interpret \(g_\Omega(\cdot,w)\) as the negative pluricomplex Green function with a logarithmic pole at \(w\):
\[
 g_\Omega(z,w)=\left(\sup\left\{u(z):u\in\operatorname{PSH}(\Omega),\ u\leq0,\ \limsup_{\zeta\to w}\bigl(u(\zeta)-\log\|\zeta-w\|\bigr)<+\infty\right\}\right)^*.
\]
The star denotes upper semicontinuous regularization in \(z\), and \(\|\cdot\|\) is the Euclidean norm. The supplied excerpt does not explicitly specify the ambient dimension or define \(g_\Omega\); the formulation here uses these standard pluripotential-theoretic conventions, which require source confirmation.

The source gives the fixed-pole boundary property as background: for a hyperconvex domain and each fixed \(w\in\Omega\),
\[
 \lim_{z\to\partial\Omega}g_\Omega(z,w)=0.
\]
The question attributed in the record to Blocki is the following. For every bounded hyperconvex domain \(\Omega\subset\C^n\), is it true that
\[
 g_\Omega(\cdot,w)\longrightarrow0\quad\text{locally uniformly on }\Omega
 \quad\text{as }w\to\partial\Omega\text{ within }\Omega?
\]
Precisely, for every nonempty compact set \(K\subset\Omega\) and every \(\varepsilon>0\), must there exist
\[
 0<\delta<\operatorname{dist}(K,\partial\Omega)
\]
such that, for every \(w\in\Omega\) satisfying \(\operatorname{dist}(w,\partial\Omega)<\delta\), one has
\[
 \sup_{z\in K}|g_\Omega(z,w)|<\varepsilon?
\]
The restriction on \(\delta\) ensures that the pole lies outside \(K\). No boundary smoothness assumption is supplied.
\subsection{Short English statement}
For a bounded hyperconvex domain, does its negative pluricomplex Green function tend uniformly to zero on every compact subset as the pole approaches any part of the boundary? The supplied record attributes the question to Blocki and labels it open; that status has not been independently verified.
\subsection{Sources}
\begin{itemize}
\item Local UnsolvedMath registry: permanent ID \(800895\), source problem code \texttt{catalogue entry}, title ``Open Problems in Pluripotential Theory --- Question 26 (Blocki).'' Supplied upstream revision: \texttt{d371ee917ac9647fbb06cd6fc529deec6f3dda82}.
\item Dinew, Guedj, Zeriahi, \emph{Open Problems in Pluripotential Theory} (2015), Question 26. Supplied source URL: \url{https://arxiv.org/abs/1511.00705}. The registry lists Dinew, Guedj, Zeriahi as proposers and records attribution/context as Blocki.
\item The supplied provenance records access on 2026-10-16 and extraction by \texttt{arxiv\_latex}. This is supplied provenance, not a claim of independent source access or verification. The registry labels the problem open; no independent literature assessment is asserted.
\end{itemize}
\end{document}
